\subsection{\texorpdfstring{ALGEBRAS OF TYPE \(C_{l}\) ( \(l \geq 1\) )}{ALGEBRAS OF TYPE C\_\{l\} ( l \textbackslash geq 1 )}}

\begin{enumerate}
\def\labelenumi{(\Roman{enumi})}
\tightlist
\item
  Let \(\varPsi\) be a non-degenerate alternating bilinear form on a vector space \(V\) of finite dimension \(2l \geq 2\) ; the set of endomorphisms \(x\) of \(V\) such that \(\varPsi(xv, v') + \varPsi(v, xv') = 0\) for all \(v, v' \in V\) is a semi-simple Lie subalgebra of \(\mathfrak{sl}(V)\) (Chap. I, §6, no. 7, Prop. 9). We denote it by \(\mathfrak{sp}(\varPsi)\) and call it the symplectic Lie algebra associated to \(\varPsi\) .
\end{enumerate}

By Algebra, Cap. IX, §4, no. 2, V can be written as the direct sum of two maximal totally isotropic subspaces F and \(F'\) , which are in duality relative to \(\Psi\) . Let \((e_{i})_{1\leq i\leq l}\) be a basis of F, and \((e_{-i})_{1\leq i\leq l}\) the dual basis of \(F'\) . Then

\[
(e _ {1}, \dots , e _ {l}, e _ {- l}, \dots , e _ {- 1})
\]

is a basis of V; we say that it is a Witt basis (or symplectic basis) of V. The matrix of \(\varPsi\) with respect to this basis is the square matrix of order 2l

\[
J = \left( \begin{array}{c c} 0 & s \\ - s & 0 \end{array} \right)
\]

where s is the square matrix of order l all of whose entries are zero except those on the second diagonal which are equal to 1, cf.~no. 2.I.

The algebra \(\mathfrak{g} = \mathfrak{sp}(\varPsi)\) can be identified with the algebra \(\mathfrak{sp}(2l,k)\) of square matrices \(a\) of order \(2l\) such that \(a = -J^{-1t}aJ = J^t aJ\) (Algebra, Chap. IX, §1, no. 10, formulas (50)), that is of the form

\[
a = \left( \begin{array}{c c} A & B \\ C & - s ^ {t} A s \end{array} \right)
\]

where A, B, C are square matrices of order l such that \(B = s^{t}Bs\) and \(C = c^{t}Cs\) ; in other words, B and C are symmetric with respect to the second diagonal. It follows that

\[
\dim \mathfrak {g} = l ^ {2} + 2 \frac {l (l + 1)}{2} = l (2 l + 1).
\]

Let \(\mathfrak{h}\) be the set of diagonal matrices in \(\mathfrak{g}\) . This is a commutative subalgebra of \(\mathfrak{g}\) , with basis the elements \(H_{i} = E_{i,i} - E_{-i,-i}\) for \(1 \leq i \leq l\) . Let \((\varepsilon_{i})_{1 \leq i \leq l}\) be the dual basis of \((H_{i})\) . For \(1 \leq i < j \leq l\) , put

\[
\left\{ \begin{array}{l l} X _ {2 \varepsilon_ {i}} & = E _ {i, - i} \\ X _ {- 2 \varepsilon_ {i}} & = - E _ {- i, i} \\ X _ {\varepsilon_ {i} - \varepsilon_ {j}} & = E _ {i, j} - E _ {- j, - i} \\ X _ {- \varepsilon_ {i} + \varepsilon_ {j}} & = - E _ {j, i} + E _ {- i, - j} \\ X _ {\varepsilon_ {i} + \varepsilon_ {j}} & = E _ {i, - j} + E _ {j, - i} \\ X _ {- \varepsilon_ {i} - \varepsilon_ {j}} & = - E _ {- i, j} - E _ {- j, i}. \end{array} \right.\tag{6}
\]

It is easily verified that these elements form a basis of a complement of \(\mathfrak{h}\) in \(\mathfrak{g}\) and that, for \(h\in \mathfrak{h}\) ,

\[
[ h, X _ {\alpha} ] = \alpha (h) X _ {\alpha}\tag{7}
\]

for all \(\alpha \in \mathbb{R}\) , where \(\mathbb{R}\) is the set of the \(\pm 2\varepsilon_{i}\) and the \(\pm \varepsilon_{i} \pm \varepsilon_{j}\) ( \(i < j\) ). It follows that \(\mathfrak{h}\) is equal to its own normalizer in \(\mathfrak{g}\) , and hence is a Cartan subalgebra of \(\mathfrak{g}\) , that \(\mathfrak{h}\) is splitting, and that the roots of \((\mathfrak{g}, \mathfrak{h})\) are the elements of \(\mathbb{R}\) . The root system \(\mathbb{R}\) of \((\mathfrak{g}, \mathfrak{h})\) is of type \(\mathrm{C}_{l}\) for \(l \geq 2\) , and of type \(\mathrm{A}_{1}\) (in other words of type \(\mathrm{C}_{1}\) ) for \(l = 1\) (Chap. VI, §4, no. 6.I extended to the case \(l = 1\) ). Consequently, \(\mathfrak{g}\) is a splittable simple Lie algebra of type \(\mathrm{C}_{l}\) .

Every splitting Cartan subalgebra of g is transformed into h by an elementary automorphism, hence by an element of the symplectic group \(\mathbf{Sp}(\varPsi)\) (cf.~(VII)), and consequently is the set \(h_{\beta}\) of elements of g whose matrix with respect to a Witt basis \(\beta\) of g is diagonal. It is immediately verified that the only vector subspaces of V stable under \(h_{\beta}\) are those generated by a subset of \(\beta\) .

We have \(\mathfrak{sp}(2,k) = \mathfrak{sl}(2,k)\) . On the other hand, the algebras \(\mathfrak{sp}(4,k)\) and \(\mathfrak{o}_S(5,k)\) have the same root system, and hence are isomorphic (cf.~Exerc. 3). From now on, we assume that \(l \geq 2\) .

\begin{enumerate}
\def\labelenumi{(\Roman{enumi})}
\setcounter{enumi}{1}
\tightlist
\item
  The root system \(R^{\vee}\) is determined by means of Chap. VI, §4, no. 6.I and 6.V; we find that
\end{enumerate}

\[
H _ {2 \varepsilon_ {i}} = H _ {i}, H _ {\varepsilon_ {i} - \varepsilon_ {j}} = H _ {i} - H _ {j}, H _ {\varepsilon_ {i} + \varepsilon_ {j}} = H _ {i} + H _ {j}.
\]

\begin{enumerate}
\def\labelenumi{(\Roman{enumi})}
\setcounter{enumi}{2}
\tightlist
\item
  Put \(\alpha_{1} = \varepsilon_{1} - \varepsilon_{2},\ldots ,\alpha_{l - 1} = \varepsilon_{l - 1} - \varepsilon_{l},\alpha_{l} = 2\varepsilon_{l}\) . By Chap. VI, §4, no. 6.II, \(\{\alpha_1,\dots ,\alpha_l\}\) is a basis B of R; the positive roots relative to B are the \(2\varepsilon_{i}\) and the \(\varepsilon_{i}\pm \varepsilon_{j}\) ( \(i < j\) ). The corresponding Borel subalgebra \(\mathfrak{b}\) is the set of upper triangular matrices in \(\mathfrak{g}\) .
\end{enumerate}

Let \(\delta\) be an isotropic flag in V (that is, whose elements are all totally isotropic subspaces for \(\varPsi\) ), and let \(\mathfrak{p}_{\delta}\) be the subalgebra consisting of the elements of \(\mathfrak{g}\) leaving stable the elements of \(\delta\) . We show as in no. 2.III that the map \(\delta \mapsto \mathfrak{p}_{\delta}\) is a bijection from the set of isotropic flags (resp. the maximum isotropic flags) to the set of parabolic (resp. Borel) subalgebras of \(\mathfrak{g}\) ; we have \(\mathfrak{p}_{\delta} \supset \mathfrak{p}_{\delta'}\) if and only if \(\delta \subset \delta'\) .

\begin{enumerate}
\def\labelenumi{(\Roman{enumi})}
\setcounter{enumi}{3}
\tightlist
\item
  The fundamental weights corresponding to \(\alpha_{1},\ldots,\alpha_{l}\) are, by Chap. VI, §4, no. 6.VI, the \(\varpi_{i}=\varepsilon_{1}+\cdots+\varepsilon_{i}\ (1\leq i\leq l)\) .
\end{enumerate}

We are going to show how the fundamental representation \(\sigma_{r}\) of weight \(\varpi_{r}\) can be realised as a subrepresentation of \(\bigwedge^{r}\sigma\) , where \(\sigma\) is the identity representation of g on V, and for this we shall study the decomposition of the representation \(\bigwedge\sigma\) of g on the exterior algebra \(\bigwedge V\) .

Let \((e_{i}^{*})\) be the basis of \(V^{*}\) dual to \((e_{i})\) . The alternating bilinear form \(\varPsi\) can be identified with an element \(\varGamma^{*}\in\bigwedge^{2}V^{*}\) (Algebra, Chap. III, §7, no. 4, Prop. 7 and §11, no. 10) and it is easy to verify that

\[
\Gamma^ {*} = - \sum_ {i = 1} ^ {l} e _ {i} ^ {*} \wedge e _ {- i} ^ {*}.
\]

Let \(\varPsi^{*}\) be the inverse form of \(\varPsi\) (Algebra, Chap. IX, §1, no. 7); it is immediate that

\[
\varPsi^ {*} (e _ {i} ^ {*}, e _ {j} ^ {*}) = 0
\]

for \(i \neq -j\) and \(\Psi^{*}(e_{i}^{*}, e_{-i}^{*}) = -1\) for \(1 \leq i \leq l\) . If we identify \(\Psi^{*}\) with an element \(\Gamma \in \bigwedge^{2}\mathrm{V}\) , then

\[
\Gamma = \sum_ {i = 1} ^ {l} e _ {i} \wedge e _ {- i}.
\]

Denote by \(X_{-}\) the endomorphism of \(\bigwedge V\) given by the left exterior product with \(\Gamma\) and by \(X_{+}\) the endomorphism of \(\bigwedge V\) given by the left interior product with \(-\Gamma^{*}\) :

\[
X _ {-} u = \left(\sum_ {i = 1} ^ {l} e _ {i} \wedge e _ {- i}\right) \wedge u,
\]

\[
X _ {+} u = \left(\sum_ {i = 1} ^ {l} e _ {i} ^ {*} \wedge e _ {- i} ^ {*}\right) \lrcorner u.
\]

To calculate \(X_{+}\) and \(X_{-}\) , introduce a basis of \(\bigwedge V\) in the following way: for any triplet (A, B, C) formed by three disjoint subsets of \([1, l]\) , put

\[
e _ {\mathrm{A,B,C}} = e _ {a _ {1}} \wedge \dots \wedge e _ {a _ {m}} \wedge e _ {- b _ {1}} \wedge \dots \wedge e _ {- b _ {n}} \wedge e _ {c _ {1}} \wedge e _ {- c _ {1}} \wedge \dots \wedge e _ {c _ {p}} \wedge e _ {- c _ {p}}
\]

where \((a_{1},\ldots,a_{m})\) (resp. \((b_{1},\ldots,b_{n}),(c_{1},\ldots,c_{p})\) ) are the elements of A (resp. B, C) arranged in increasing order. We obtain in this way a basis of \(\bigwedge V\) and simple calculations show that

\[
X _ {-}. e _ {\mathrm{A,B,C}} = \sum_ {j \in \{1, l \}, j \notin \mathrm{A} \cup \mathrm{B} \cup \mathrm{C}} e _ {\mathrm{A,B,C} \cup \{j \}}\tag{8}
\]

\[
X _ {+}. e _ {\mathrm{A,B,C}} = - \sum_ {j \in \mathrm{C}} e _ {\mathrm{A,B,C-} \{j \}}.\tag{9}
\]

Let H be the endomorphism of \(\bigwedge V\) that reduces to multiplication by \((l - r)\) on \(\bigwedge^{r}V\) ( \(0 \leq r \leq 2l\) ). It is painless to verify (cf.~Exerc. 19) that

\[
\begin{array}{c} {[ X _ {+}, X _ {-} ] = - H} \\ {[ H, X _ {+} ] = 2 X _ {+}} \\ {[ H, X _ {-} ] = - 2 X _ {-}.} \end{array}
\]

In other words, the vector subspace \(\mathfrak{s}\) generated by \(X_{+}, X_{-}\) and \(H\) is a Lie subalgebra of \(\operatorname{End}(\bigwedge \mathrm{V})\) , isomorphic to \(\mathfrak{sl}(2, k)\) , and \(\bigwedge^{r} \mathrm{V}\) is the subspace of elements of weight \(l - r\) . Denote by \(\mathrm{E}_{r}\) the subspace of \(\bigwedge^{r} \mathrm{V}\) consisting of the primitive elements, that is, \(\mathrm{E}_{r} = (\bigwedge^{r} \mathrm{V}) \cap \mathrm{Ker} X_{+}\) . It follows from §1 that, for \(r < l\) , the restriction of \(X_{-}\) to \(\bigwedge^{r} \mathrm{V}\) is injective and that, for \(r \leq l\) , \(\bigwedge^{r} \mathrm{V}\) decomposes as a direct sum

\[
\begin{array}{c} \bigwedge^ {r} \mathrm{V} = \mathrm{E} _ {r} \oplus X _ {-} (\mathrm{E} _ {r - 2}) \oplus X _ {-} ^ {2} (\mathrm{E} _ {r - 4}) \oplus \dots \\ = \mathrm{E} _ {r} \oplus X _ {-} (\bigwedge^ {r - 2} \mathrm{V}). \end{array}
\]

This shows in particular that \(\dim \mathrm{E}_r = \binom{2l}{r} - \binom{2l}{r-2}\) for \(0 \leq r \leq l\) .

On the other hand, the very definition of \(\mathfrak{sp}(\varPsi)\) shows that \(\Gamma^{*}\) is annihilated by the second exterior power of the dual of \(\sigma\) . Similarly, \(\Gamma\) is annihilated by \(\bigwedge^{2}\sigma\) . We deduce immediately that \(X_{+}\) and \(X_{-}\) , and hence also \(H\) , commute with the endomorphisms \(\bigwedge\sigma(g)\) for \(g\in\mathfrak{g}\) . Consequently, the subspaces \(\mathrm{E}_{r}\) for \(0\leq r\leq l\) are stable under \(\bigwedge^{r}\sigma\) ; we shall show that the restriction of \(\bigwedge^{r}\sigma\) to \(\mathrm{E}_{r}\) is a fundamental representation \(\sigma_{r}\) of weight \(\varpi_{r}\) ( \(1\leq r\leq l\) ).

We remark first of all that the weights of \(\bigwedge^{r}\sigma\) relative to \(\mathfrak{h}\) are the

\[
\varepsilon_ {i _ {1}} + \dots + \varepsilon_ {i _ {k}} - (\varepsilon_ {j _ {1}} + \dots + \varepsilon_ {j _ {r - k}}),
\]

where \(i_1, \ldots, i_k\) (resp. \(j_1, \ldots, j_{r-k}\) ) are distinct elements of \([1, l]\) ; thus, the highest weight of \(\bigwedge^r \sigma\) is indeed

\[
\varpi_ {r} = \varepsilon_ {1} + \dots + \varepsilon_ {r}
\]

and the vectors of weight \(\varpi_{r}\) are those proportional to \(e_{1} \wedge \cdots \wedge e_{r} = e_{\{1,\ldots,r\},\varnothing,\varnothing}\) . Formula (9) shows that \(e_{1} \wedge \cdots \wedge e_{r} \in E_{r}\) . Thus, it suffices to prove that the restriction of \(\bigwedge^{r}\sigma\) to \(E_{r}\) is irreducible.

If \(s \in \mathbf{Sp}(\varPsi)\) , the extension of s to \(\bigwedge V\) (resp. \(\bigwedge V^{*}\) ) fixes \(\varGamma\) (resp. \(\varGamma^{*}\) ), and hence commutes with \(X_{+}\) and \(X_{-}\) and leaves \(E_{r}\) stable. Consequently, \(E_{r}\) contains the vector subspace \(F_{r}\) generated by the transforms of \(e_{1} \wedge \cdots \wedge e_{r}\) by \(\mathbf{Sp}(\varPsi)\) . The theorem of Witt shows that these are the non-zero decomposable r-vectors such that corresponding vector subspace of V is a totally isotropic subspace, r-vectors that we shall call isotropic.

Lemma 2. For \(1 \leq r \leq l\) , let \(\mathrm{F}_r\) be the subspace of \(\bigwedge^r\mathrm{V}\) generated by the isotropic \(r\) -vectors. Then

\[
\bigwedge^ {r} \mathrm{V} = \mathrm{F} _ {r} + X _ {-} (\bigwedge^ {r - 2} \mathrm{V}) = \mathrm{F} _ {r} + \left(\sum_ {i = 1} ^ {l} e _ {i} \wedge e _ {- i}\right) \wedge \bigwedge^ {r - 2} \mathrm{V}.
\]

We show first of all how Lemma 2 implies our assertion. Since \(F_{r} \subset E_{r}\) and \(\mathrm{E}_{r} \cap X_{-}(\bigwedge^{r-2}\mathrm{V}) = \{0\}\) , the lemma implies that \(F_{r} = E_{r}\) . On the other hand, let \(s \in \mathbf{Sp}(\varPsi)\) ; the automorphism \(a \mapsto sas^{-1}\) of \(\operatorname{End}(V)\) preserves g and induces on it an element of \(\operatorname{Aut}_{0}(\mathfrak{g})\) (cf.~(VII)), and hence transforms every irreducible representation of g into an equivalent representation (§7, no. 1, Prop. 2). Since \(e_{1} \wedge \cdots \wedge e_{r}\) belongs to an irreducible component of \(\bigwedge^{r}\sigma\) , and since \(E_{r} = F_{r}\) is generated by the transforms of \(e_{1} \wedge \cdots \wedge e_{r}\) by \(\mathbf{Sp}(\varPsi)\) , it follows that the representation of g on \(E_{r}\) is isotypical. But the multiplicity of its highest weight \(\varpi_{r}\) is 1, so it is irreducible.

It remains to prove the lemma. It is clear for r = 1. We argue by induction, and assume that \(r \geq 2\) . By the induction hypothesis, we are reduced to proving that

\[
\mathrm{F} _ {r - 1} \wedge \mathrm{V} \subset \mathrm{F} _ {r} + \Gamma \wedge \bigwedge^ {r - 2} \mathrm{V},
\]

or that, if \(y\) is a decomposable \((r - 1)\) -vector and \(x \in V\) , then

\[
z = y \wedge x \in \mathrm{F} _ {r} + \Gamma \wedge \bigwedge^ {r - 2} \mathrm{V}.
\]

Let \((f_{i})_{1\leq\pm i\leq l}\) be a Witt basis of V such that \(y=f_{1}\wedge\cdots\wedge f_{r-1}\) . It suffices to carry out the proof when \(x=f_{i}\) . If \(i\notin(1-r,-1)\) , the r-vector \(f_{1}\wedge\cdots\wedge f_{r-1}\wedge f_{i}\) is isotropic. Otherwise, we can assume, renumbering the \(f_{i}\) if necessary, that i=1-r. Then \(\Gamma=\sum_{j=1}^{l}f_{j}\wedge f_{-j}\) , so

\[
\begin{array}{l} f _ {r - 1} \wedge f _ {1 - r} = \frac {1}{l - r + 2} \left(\Gamma - \sum_ {i = 1} ^ {r - 2} f _ {i} \wedge f _ {- i} + \sum_ {j = r} ^ {l} (f _ {r - 1} \wedge f _ {1 - r} - f _ {j} \wedge f _ {- j})\right), \\ z = \frac {1}{l - r + 2} \Gamma \wedge f _ {1} \wedge \dots \wedge f _ {r - 2} \\ \qquad + \frac {1}{l - r + 2} \sum_ {j = r} ^ {l} (f _ {1} \wedge \dots \wedge f _ {r - 2}) \wedge (f _ {r - 1} \wedge f _ {1 - r} - f _ {j} \wedge f _ {- j}). \end{array}
\]

But

\[
f _ {r - 1} \wedge f _ {1 - r} - f _ {j} \wedge f _ {- j} = (f _ {r - 1} + f _ {j}) \wedge (f _ {1 - r} - f _ {- j}) - f _ {j} \wedge f _ {1 - r} + f _ {r - 1} \wedge f _ {- j}
\]

and we verify immediately that the r-vectors

\[
\begin{array}{c} f _ {1} \wedge \dots \wedge f _ {r - 2} \wedge (f _ {r - 1} + f _ {j}) \wedge (f _ {1 - r} - f _ {- j}), \\ f _ {1} \wedge \dots \wedge f _ {r - 2} \wedge f _ {j} \wedge f _ {1 - r} \text {and} f _ {1} \wedge \dots \wedge f _ {r - 1} \wedge f _ {- j} \end{array}
\]

are isotropic for \(r \leq j \leq l\) . Consequently, \(z \in F_{r} + \Gamma \wedge \bigwedge^{r-2} V\) , which completes the proof.

\begin{enumerate}
\def\labelenumi{(\Alph{enumi})}
\setcounter{enumi}{21}
\tightlist
\item
  We have \(w_{0} = -1\) , so every finite dimensional simple representation of g is orthogonal or symplectic. By Chap. VI, §4, no. 6.VI, the sum of the coordinates of \(\varpi_{r}\) with respect to \((\alpha_{1}, \ldots, \alpha_{l})\) is
\end{enumerate}

\[
1 + 2 + \dots + (r - 1) + r + r + \dots + r + \frac {r}{2}
\]

so \(\sigma_{r}\) is orthogonal for r even and symplectic for r odd.

Since \(e_{1} \wedge \cdots \wedge e_{r}\) and \(e_{-1} \wedge \cdots \wedge e_{-r}\) belong to \(E_{r}\) and since

\[
\Psi_ {(r)} \left(e _ {1} \wedge \dots \wedge e _ {r}, e _ {- 1} \wedge \dots \wedge e _ {- r}\right) = 1,
\]

we see that the restriction of \(\Psi_{(r)}\) to \(E_{r}\) is non-zero: this is, up to a constant factor, the bilinear form; it is symmetric if r is even, alternating if r is odd, and invariant under \(\sigma_{r}\) .

\begin{enumerate}
\def\labelenumi{(\Roman{enumi})}
\setcounter{enumi}{5}
\tightlist
\item
  For all \(x \in \mathfrak{g}\) , the characteristic polynomial of \(\sigma(x)\) takes the form
\end{enumerate}

\[
\mathrm{T} ^ {2 l} + f _ {1} (x) \mathrm{T} ^ {2 l - 1} + \dots + f _ {2 l} (x)
\]

where \(f_{1},\ldots ,f_{2l}\) are invariant polynomial functions on \(\mathfrak{g}\) .

If \(x = \xi_{1}H_{1} + \cdots + \xi_{l}H_{l} \in \mathfrak{h}\) , the \(f_{i}(x)\) are, up to sign, the elementary symmetric functions of \(\xi_{1}, \ldots, \xi_{l}, -\xi_{1}, \ldots, -\xi_{l}\) ; these symmetric functions are zero in odd degrees, and

\[
\mathrm{T} ^ {2 l} + f _ {2} (x) \mathrm{T} ^ {2 l - 2} + \dots + f _ {2 l} (x) = (\mathrm{T} ^ {2} - \xi_ {1} ^ {2}) \ldots (\mathrm{T} ^ {2} - \xi_ {l} ^ {2}).
\]

As in no. 2.VI, it follows that \(f_{1}=f_{3}=f_{5}=\cdots=0\) , and that

\[
(f _ {2}, f _ {4}, \dots , f _ {2 l})
\]

is an algebraically free family generating the algebra of invariant polynomial functions on \(\mathfrak{g}\) .

\begin{enumerate}
\def\labelenumi{(\Roman{enumi})}
\setcounter{enumi}{6}
\tightlist
\item
  Since the only automorphism of the Dynkin graph is the identity, we have \(\mathrm{Aut}(\mathfrak{g}) = \mathrm{Aut}_0(\mathfrak{g})\) .
\end{enumerate}

Let \(\Sigma\) be the group of similarities of \(V\) relative to \(\Psi\) (Algebra, Chap. IX, §6, end of no. 5). One proves as in no. 2.VII that the automorphisms of \(\mathfrak{g}\) are the maps \(x \mapsto sxs^{-1}\) where \(s \in \Sigma\) , so that \(\operatorname{Aut}(\mathfrak{g}) = \operatorname{Aut}_0(\mathfrak{g})\) can be identified with \(\Sigma / k^*\) .

For all \(s \in \Sigma\) , let \(\mu(s)\) be the multiplier of s. The map \(s \mapsto \mu(s) \mod k^{*2}\) from \(\Sigma\) to \(k^{*}/k^{*2}\) is a homomorphism whose kernel contains \(k^{*}.1\) , and hence gives a homomorphism \(\lambda\) from \(\Sigma/k^{*}\) to \(k^{*}/k^{*2}\) . We have \(\mathbf{Sp}(\varPsi) \cap k^{*} = \{1, -1\}\) . Consider the sequence of homomorphisms

\[
1 \longrightarrow \mathbf {S p} (\Psi) / \{1, - 1 \} \stackrel {{\iota}} {{\longrightarrow}} \Sigma / k ^ {*} \stackrel {{\lambda}} {{\longrightarrow}} k ^ {*} / k ^ {* 2} \longrightarrow 1.\tag{10}
\]

The map \(\iota\) is injective, and \(\mathrm{Im}(\iota) \subset \mathrm{Ker}\lambda\) since the multiplier of an element of \(\mathbf{Sp}(\varPsi)\) is 1. If the multiplier of \(s \in \varSigma\) is an element of \(k^{*2}\) , there exists \(\nu \in k^*\) such that \(\nu s \in \mathbf{Sp}(\varPsi)\) ; thus, \(\mathrm{Im}(\iota) = \mathrm{Ker}(\lambda)\) . In summary, the sequence (10) is exact. We identify \(\mathbf{Sp}(\varPsi)/\{1, -1\}\) with a subgroup of \(\varSigma/k^*\) . Since \(k^*/k^{*2}\) is commutative, \(\mathbf{Sp}(\varPsi)/\{1, -1\}\) contains the derived group of \(\varSigma/k^*\) . Thus, \(\mathrm{Aut}_e(\mathfrak{g})\) is contained in \(\mathbf{Sp}(\varPsi)/\{1, -1\}\) ( \(\S\) 11, no. 2, Prop. 3). In fact, it is equal to it, and \(\mathrm{Aut}(\mathfrak{g})/\mathrm{Aut}_e(\mathfrak{g})\) is identified with \(k^*/k^{*2}\) (Exerc. 9).

\begin{enumerate}
\def\labelenumi{(\Roman{enumi})}
\setcounter{enumi}{7}
\tightlist
\item
  The canonical bilinear form \(\Phi_{\mathrm{R}}\) on \(\mathfrak{h}^*\) is given by
\end{enumerate}

\[
\varPhi_ {\mathrm{R}} (\xi_ {1} \varepsilon_ {1} + \dots + \xi_ {l} \varepsilon_ {l}, \xi_ {1} ^ {\prime} \varepsilon_ {1} + \dots + \xi_ {l} ^ {\prime} \varepsilon_ {l}) = \frac {1}{4 (l + 1)} (\xi_ {1} \xi_ {1} ^ {\prime} + \dots + \xi_ {l} \xi_ {l} ^ {\prime})
\]

(Chap. VI, §4, no. 6.V). Thus, the inverse form of \(\varPhi_{\mathrm{R}}\) , that is, the restriction to \(\mathfrak{h}\) of the Killing form, is

\[
\Phi (\xi_ {1} H _ {1} + \dots + \xi_ {l} H _ {l}, \xi_ {1} ^ {\prime} H _ {1} + \dots + \xi_ {l} ^ {\prime} H _ {l}) = 4 (l + 1) (\xi_ {1} \xi_ {1} ^ {\prime} + \dots + \xi_ {l} \xi_ {l} ^ {\prime}).
\]

\begin{enumerate}
\def\labelenumi{(\Roman{enumi})}
\setcounter{enumi}{8}
\tightlist
\item
  Recall the \(X_{\alpha}\) defined by formulas (6) ( \(\alpha \in R\) ). It is easily verified that \([X_{\alpha}, X_{-\alpha}] = -H_{\alpha}\) for \(\alpha \in R\) . On the other hand, the map \(\theta : a \mapsto -^{t}a\) is an automorphism of g and \(\theta(X_{\alpha}) = X_{-\alpha}\) for all \(\alpha \in R\) . Consequently, \((X_{\alpha})_{\alpha \in \mathbb{R}}\) is a Chevalley system in \((\mathfrak{g}, \mathfrak{h})\) .
\end{enumerate}

Assume that \(k = \mathbf{Q}\) . The Cartan subalgebra \(\mathfrak{h}\) has two permissible lattices \(\mathrm{Q}(\mathrm{R}^{\vee}) = \sum_{i=1}^{l} \mathbf{Z}.H_i\) and \(\mathrm{P}(\mathrm{R}^{\vee}) = \mathrm{Q}(\mathrm{R}^{\vee}) + \frac{1}{2}\mathbf{Z}. \sum_{i=1}^{l} H_i\) (Chap. VI, §4, no. 5.VIII). We see that \(\mathrm{Q}(\mathrm{R}^{\vee})\) is the set of matrices with integer entries belonging to \(\mathfrak{h}\) . It follows that the Chevalley order \(\mathrm{Q}(\mathrm{R}^{\vee}) + \sum_{\alpha \in \mathrm{R}} \mathbf{Z}.X_{\alpha}\) is the set \(\mathfrak{sp}(2l,\mathbf{Z})\) of matrices in \(\mathfrak{g}\) with integer entries.

Consider the reductive Lie algebra \(\mathfrak{sp}(\varPsi) + \mathbf{Q}.1\) . It is easy to see that the set of its elements with integer entries is a Chevalley order, whose projection onto \(\mathfrak{sp}(\varPsi)\) parallel to \(\mathbf{Q}.1\) is the Chevalley order \(\mathrm{P}(\mathrm{R}^{\vee}) + \sum \mathbf{Z}.X_{\alpha}\) .

Finally, \(X_{\alpha}^{2} = 0\) for all \(\alpha \in \mathbb{R}\) . It follows that the lattice \(\overline{\mathcal{V}}\) in \(\mathrm{V}\) generated by the \(e_i\) is admissible for the Chevalley order \(\mathfrak{sp}(2l,\mathbf{Z})\) . The same holds for the lattice \(\mathrm{E}_r \cap \bigwedge^r \mathcal{V}\) in \(\mathrm{E}_r\) .

Finally, \(E_{r}\) has an admissible lattice for the Chevalley order

\[
\mathrm{P} (\mathrm{R} ^ {\vee}) + \sum \mathbf {Z}. X _ {\alpha}
\]

only if \(r\) is even; then \(\mathrm{E}_r\cap \bigwedge^r\mathcal{V}\) is such a lattice.

